\chapter{Field Protocols}
\label{app:protocols}

\section{Shoot census}
\label{sec:protocol-census}

Each quadrat was located by its fixed corner pegs and photographed from
above before counting. Divers counted every vegetative and generative
shoot rooted inside the frame; shoots on the boundary were counted on two
sides only. Counts were repeated by a second diver in one quadrat in five,
and the census was repeated whenever the two counts differed by more than
ten per cent.

\section{Sediment coring and dating}
\label{sec:protocol-cores}

Cores were extruded on the day of collection and sliced at two-centimetre
intervals. Each slice was freeze-dried, homogenised and split for
organic-carbon and radionuclide analyses. \Cref{tab:protocol} summarises
the quality controls applied to every batch.

\begin{table}[hb]
  \centering\small
  \caption{Quality controls for sediment analyses.}
  \label{tab:protocol}
  \begin{tabular*}{\linewidth}{@{\extracolsep{\fill}}l l r@{}}
    \toprule
    Analysis & Control & Tolerance \\
    \midrule
    Organic carbon & Certified reference sediment & \(\pm 3\,\%\) \\
                   & Duplicate slice, one in ten   & \(\pm 5\,\%\) \\
    Lead-210       & Blank counting, weekly        & \(< 0.2\) Bq \\
                   & Spiked standard, per batch    & \(\pm 4\,\%\) \\
    \bottomrule
  \end{tabular*}
\end{table}

Ages were derived with the constant-rate-of-supply model, rejecting any
core in which the excess lead-210 profile was non-monotonic in its upper
twenty centimetres \parencite{tetley2020dating}.
